% Generated draft: requires mathematical and source review.
\documentclass[11pt]{article}
\usepackage[margin=25mm]{geometry}
\usepackage{fontspec}
\setmainfont{Latin Modern Roman}
\usepackage{amsmath,amssymb,mathtools,unicode-math}
\setmathfont{Latin Modern Math}
\usepackage{enumitem,needspace,xurl,hyperref}
\setlength{\parindent}{0pt}
\setlength{\parskip}{5pt}
\setlength{\emergencystretch}{3em}
\newcommand{\N}{\mathbb{N}}
\newcommand{\Z}{\mathbb{Z}}
\newcommand{\Q}{\mathbb{Q}}
\newcommand{\R}{\mathbb{R}}
\newcommand{\C}{\mathbb{C}}
\newcommand{\F}{\mathbb{F}}
\newcommand{\A}{\mathbb{A}}
\newcommand{\PP}{\mathbb{P}}
\newcommand{\E}{\mathbb{E}}
\newcommand{\eps}{\varepsilon}
\newcommand{\dd}{\,\mathrm{d}}
\begin{document}
\section*{UnsolvedMath ID: 800014}
\textbf{Source problem code:} catalogue entry
\subsection{Definitions and mathematical statement}
A commutative semigroup is a nonempty set \(S\) with a binary operation satisfying \((xy)z=x(yz)\) and \(xy=yx\) for all \(x,y,z\in S\). No identity element is required. A variety of commutative semigroups is a class of such semigroups axiomatized by identities in the signature consisting of one binary operation.

The supplied question, attributed to S. Kublanovsky, is: for every variety \(\mathcal V\) of commutative semigroups, is its ``endomorphic function'' ``identical''?

The source additionally states that the variety of all commutative semigroups is ``endomorphic'', referring to [kubl]. These technical expressions are retained without interpretation: the supplied record does not define an endomorphic function, its domain or codomain, the meaning of ``identical'', or the property of being an endomorphic variety. Consequently, the intended predicate on \(\mathcal V\) cannot be specified precisely from this record alone. In particular, interpreting ``identical'' as equality to an identity map would require additional source evidence.
\subsection{Short English statement}
S. Kublanovsky asks whether every variety of commutative semigroups has an ``identical'' endomorphic function. The supplied record leaves both technical expressions undefined, so the intended mathematical assertion requires clarification from the source and its reference [kubl].
\subsection{Sources}
\begin{itemize}
\item Local UnsolvedMath registry, supplied record with permanent ID \(800014\), source problem code \texttt{catalogue entry}, titled \emph{The Lyapin's Notebook --- Problem 14}. Its supplied status is open; that status has not been independently verified.
\item Bershadsky, Kublanovsky, Mashevitzky, \emph{The Lyapin's Notebook: Unsolved Problems in Semigroup Theory} (2026), Problem 14. Supplied source URL: \url{https://arxiv.org/abs/2604.04763}. The record reports access on 2026-10-16 and extraction by \texttt{arxiv\_latex}; these are supplied provenance details, not an independent source visit.
\item The individual question is attributed in the statement to S. Kublanovsky. The registry lists Bershadsky, Kublanovsky, and Mashevitzky as proposers and gives the year 2026.
\item The statement cites [kubl] for the assertion that the variety of commutative semigroups is endomorphic. No bibliographic entry or URL for [kubl] is supplied, so this reference and its intended definitions remain unresolved.
\end{itemize}
\end{document}
